\begin{pastexam}[label={px:ntu-113-dsa-03}]{臺大資工所 113 學年度 第 3 題}
（單選；答錯倒扣 1 分）

(5 points) A Dynamic Array (abbreviated as DArray) is a type of array that permits the insertion and
deletion of elements at the end to dynamically adjust its size. DArray $D$ maintains three crucial
attributes, including the underlying array $D$.data, the logical data size $D$.size, and the actual
array capacity $D$.capacity. The invariant $D$.size $\le$ $D$.capacity holds at any given moment.
The Insert($D$, $x$) operation appends the element $x$ to the end of $D$.data. In case $D$.size $=$
$D$.capacity, it triggers a call to Resize($D$) before appending $x$, where Resize($D$) allocates a
new and larger array, moves all the data to it, increases the value of $D$.capacity, sets up
$D$.data, and finally frees the old array. Two implementations of Resize($D$) are provided:
ResizeA($D$) increases $D$.capacity by a constant (e.g., 10) while ResizeB($D$) doubles
$D$.capacity each time.

To initialize DArray $D$, we set $D$.size to 0, $D$.capacity to 1, and allocate a size 1 array to
$D$.data. Assuming that both allocating and freeing arrays of length $n$ take $O(n)$ time, and moving
a single element takes $O(1)$ time, please select the correct description(s) below.
\begin{enumerate}[label=\Alph*.]
\item DArray is not suitable for random access because its memory may be relocated.
\item DArray can only be allocated in the heap memory due to its dynamic memory nature.
\item If ResizeA($D$) is considered, the amortized time complexity of Insert($D$, $x$) is $O(1)$.
\item If ResizeB($D$) is considered, the amortized time complexity of Insert($D$, $x$) is $O(1)$.
\item In the Delete($D$) operation, if $D$.size $>0$, $D$.size is decreased by one. If, at this
stage, $D$.size $<\lceil D.\text{capacity}/2\rceil$, we create a new array $d'$ with a size of
$\lceil D.\text{capacity}/2\rceil$, transfer the data to $d'$, set up $d'$ as the new data array for
$D$.data, and subsequently free the old array. Then, the amortized time complexity of mixed Insert
and Delete operations can be $O(1)$, where Resize($D$) can be either ResizeA($D$) or ResizeB($D$).
\end{enumerate}
\end{pastexam}

\begin{pastsoln}
答案是 D。

這是資結書第~3.4~節動態陣列的內容。 Resize 一次的成本是「配置新陣列 $+$ 搬 $\mathit{size}$ 個
元素 $+$ 釋放舊陣列」，都和容量同一級，以下記成 $\Theta(\mathit{capacity})$。
\begin{itemize}
\item A 錯。 動態陣列的元素仍然放在一塊連續的陣列裡，取第 $i$ 個是 $D.\mathit{data}[i]$，
$O(1)$。 搬家只是換了 $D.\mathit{data}$ 指到的位置，透過 $D.\mathit{data}$ 用索引存取不受影響；
會失效的只是「事先記下來、直接指向舊陣列某一格的指標」。 C++ 的 \texttt{vector} 就是這種結構，
它的賣點正是隨機存取。
\item B 錯。 題目只要求 Resize 能「配置一塊新陣列、釋放舊陣列」，這可以由任何配置器提供：
一般用 \texttt{malloc} 從堆積（heap）拿，但也可以從事先保留好的靜態記憶體池切一塊，嵌入式系統
常這樣做。 何況 $D$ 本身（三個欄位）可以是函式裡的區域變數，放在堆疊上。 「只能放在 heap」
這個「只能」不成立。
\item C 錯。 每次只加 $10$，容量依序是 $1, 11, 21, 31, \ldots$。 連續插入 $n$ 個時，第
$t$ 次搬家發生在 $\mathit{size} = 10t - 9$ 的時候（$t = 1, 2, \ldots$），大約有 $n/10$ 次，成本合計
\[
\sum_{t=1}^{\approx n/10} \Theta(10t) = \Theta\Bigl(10 \cdot \frac{(n/10)^2}{2}\Bigr) = \Theta(n^2),
\]
平攤到 $n$ 次插入，每次 $\Theta(n)$，不是 $O(1)$（資結書第~3.4~節「為什麼要乘以倍數，不能加固定
的量」）。
\item D 對。 容量依序是 $1, 2, 4, \ldots$，連續插入 $n$ 個時搬家發生在 $\mathit{size} = 1, 2, 4,
\ldots, 2^{\lfloor \lg (n-1) \rfloor}$，成本合計 $\Theta(1 + 2 + 4 + \cdots + 2^{\lfloor \lg (n-1)
\rfloor}) = \Theta(n)$（等比級數，最多是最後一項的兩倍）。 加上 $n$ 次放入元素的 $\Theta(n)$，總成
本 $\Theta(n)$，平攤 $O(1)$（聚集分析，資結書第~3.4.1~節）。
\item E 錯。 「少於一半就減半」會在邊界來回搬家，這正是資結書第~3.4.2~節說「要縮到 $1/4$ 才減半」
的原因。 用 ResizeB 舉例：設容量 $m$（$m \ge 2$ 是 $2$ 的次方）而且放滿 $m$ 個。
  \begin{enumerate}[label=(\arabic*)]
  \item Insert：放滿了，先加倍成 $2m$、搬 $m$ 個，之後 $\mathit{size} = m + 1$。
  \item Delete：$\mathit{size} = m$，$m < \lceil 2m/2 \rceil = m$ 不成立，不縮。
  \item Delete：$\mathit{size} = m - 1 < m$，縮成容量 $m$、搬 $m - 1$ 個。
  \item Insert：$\mathit{size} = m$，容量 $m$，又是放滿的狀態。
  \end{enumerate}
回到一開始的狀態，每 $4$ 次操作就要花 $\Theta(m)$。 先插入 $n/2$ 個讓 $m$ 和 $n$ 同一級，剩下的
操作都照這四步循環，總成本 $\Theta(n^2)$，平攤 $\Theta(n)$。 題目說「ResizeA 或 ResizeB 都可以」，
光 ResizeB 就不行；ResizeA 連只做 Insert 都不是 $O(1)$（見 C）。
\end{itemize}
本題是單選題，只有 D 對。
\end{pastsoln}
